% MyHatTileStuff/h8_rotcluster.py — Detailed Technical Explanation
% Compile with: pdflatex -shell-escape h8_rotcluster.tex
%
\documentclass[10pt,a4paper]{article}

%% ---------- Packages ----------
\usepackage[T1]{fontenc}
\usepackage[utf8]{inputenc}
\usepackage{geometry}
\geometry{margin=2cm, top=2.2cm, bottom=2.5cm}

\usepackage{microtype}
\usepackage{parskip}
\usepackage{xcolor}
\usepackage{booktabs}
\usepackage{amsmath,amssymb}
\usepackage{graphicx}
\usepackage{hyperref}
\usepackage{enumitem}
\setlist{noitemsep, topsep=2pt}
\usepackage{needspace}

%% minted — syntax highlighting (requires -shell-escape and pygments)
\usepackage{minted}
\setminted{
  fontsize=\footnotesize,
  baselinestretch=1.1,
  breaklines=true,
  linenos=false,
  frame=none,
  bgcolor=codebg
}
\definecolor{codebg}{HTML}{F5F5F5}

%% tcolorbox — side-by-side description/code panels
\usepackage[most]{tcolorbox}
\tcbuselibrary{minted,skins,breakable}

%% Colours
\definecolor{titlecolor}{HTML}{1A4D8F}
\definecolor{headercolor}{HTML}{2E6DB4}
\definecolor{rulecolor}{HTML}{AECBEF}
\definecolor{antiyellow}{HTML}{F4D03F}
\definecolor{hatblue}{HTML}{C5D6EA}
\definecolor{lvone}{HTML}{3B0F70}
\definecolor{lvtwo}{HTML}{641A80}
\definecolor{lvthree}{HTML}{A52C60}
\definecolor{lvfour}{HTML}{DE4968}
\definecolor{lvfive}{HTML}{FE9F6D}

%% Side-by-side box style
\tcbset{
  dualbox/.style={
    enhanced,
    colframe=rulecolor,
    colback=white,
    fonttitle=\bfseries\small,
    breakable,
    boxrule=0.6pt,
    arc=3pt,
    sidebyside,
    sidebyside align=top seam,
    sidebyside gap=6mm,
    lefthand ratio=0.46,
    before skip=8pt,
    after skip=8pt,
  }
}

%% Section colours
\usepackage{titlesec}
\titleformat{\section}{\large\bfseries\color{titlecolor}}{}{0em}{}[\nobreak\vspace{-2pt}\color{rulecolor}\hrule\nobreak\vspace{4pt}]
\titleformat{\subsection}{\normalsize\bfseries\color{headercolor}}{}{0em}{}

%% Keep headings attached to the content that follows them
\newcommand{\subsectionbreak}{\nobreak}
\let\polysection\section
\renewcommand\section{\needspace{0.16\textheight}\polysection}
\let\polysubsection\subsection
\renewcommand\subsection{\needspace{0.12\textheight}\polysubsection}

%% Convenient macros
\newcommand{\code}[1]{\texttt{\small #1}}
\newcommand{\file}[1]{\texttt{\footnotesize #1}}

%% ---------- Document ----------
\begin{document}

\begin{center}
  {\Huge\bfseries\color{titlecolor} h8\_rotcluster.py}\\[4pt]
  {\large\color{headercolor} H8 Rotational Symmetry: Super Cluster Assembly and Level-Coloured Overlay}\\[6pt]
  {\small \file{MyHatTileStuff/h8\_rotcluster.py} — Technical Reference}\\[2pt]
  {\small\color{gray} \today}
\end{center}

\vspace{6pt}
\noindent\rule{\linewidth}{1pt}
\vspace{4pt}

% -----------------------------------------------------------------------
\section{Purpose}

\code{h8\_rotcluster.py} bundles the H8 rotational-symmetry findings into a
single self-contained script (NumPy only).  For arrangement \#13 — the
gap-free centre of type $v[12]$ at world coordinate
$C_{13} = (-68.5,\; 0.866\ldots)$ — it performs four tasks in sequence:

\begin{enumerate}
  \item \textbf{Super cluster construction.}  Loads the base H8 super tile
    from \file{from\_python/SUPER\_SUPER\_H8\_TILE.dat} and places three
    rotated copies (copy 0/1/2) around $C_{13}$ to form a single
    3-fold symmetric ``super cluster''.  The whole assembly is rotated by
    $180°$ so that it emerges from \code{rotsym.c} the right way up.
  \item \textbf{Level-match overlay.}  For each ring level $L = 1\ldots 5$,
    reads the symmetric completions produced by \code{grow\_symmetric2}
    (\file{data/sym12/sym\_completions\_levelN.dat}) and finds the record
    whose interior best matches the super cluster.  Uses the level-5 match to
    fix a global rigid-body transform (rotation + optional reflection) from
    grow-space into super-cluster space.
  \item \textbf{SVG image.}  Renders the super cluster (anti-hats in yellow,
    hats in pale blue) with a semi-transparent level-coloured overlay and a
    bold outer boundary, written to \file{from\_python/h8\_rot\_overlay.svg}.
  \item \textbf{CSV export.}  Writes two files:
    \begin{itemize}
      \item \file{from\_python/SUPER\_SUPER\_H8\_ROT.csv} — all three copies
        one after another.
      \item \file{from\_python/LEVEL\_ROT.csv} — the grown cluster, ordered
        inside-out by level.
    \end{itemize}
\end{enumerate}

% -----------------------------------------------------------------------
\section{Background: Hat Tile Encoding}

\subsection{CSV Parameter Format}

Each hat tile in the project is stored as five numbers
$(x,\; y,\; \mathit{dir},\; \mathit{ref},\; \mathit{pt})$:

\begin{center}
\begin{tabular}{@{}lp{10cm}@{}}
\toprule
\textbf{Field} & \textbf{Meaning} \\
\midrule
$x, y$        & World-space anchor: the position in $\mathbb{R}^2$ of
                 vertex $\mathit{pt}$ of the tile. \\
$\mathit{dir}$ & Discrete orientation in the range $0\ldots 11$
                 (each step is $30°$; positive = CCW for regular hats,
                  CW for reflected). \\
$\mathit{ref}$ & Chirality flag: \texttt{False} = regular hat,
                 \texttt{True} = anti-hat (reflected). \\
$\mathit{pt}$  & Index $1\ldots 14$ of the vertex whose position
                 is given by $(x, y)$.  Acts as the ``pivot'' for the
                 local orientation frame. \\
\bottomrule
\end{tabular}
\end{center}

\noindent The five parameters are sufficient to reconstruct the full 14-vertex
polygon in world space; see \code{real\_world} below.

\subsection{Edge Lengths and Directions in \code{hv14}}

The hat tile is a 14-vertex (13-edge) polygon on the tetrille lattice.  Its
edge lengths alternate between $1$ and $\sqrt{3}$, and its turning angles at
each vertex are multiples of $30°$.

Two lookup arrays encode the canonical hat boundary:
\begin{minted}{python}
raw_len = [1,1,0,0,1,1,0,0,1,1,0,0,0,0]  # 1=sqrt(3), 0=1
raw_rot = [-1,1,4,6,3,5,8,6,9,7,10,12,12,14]
\end{minted}
\code{raw\_len[j] == 1} means edge $j$ has length $\sqrt{3}$; \code{0} means
length $1$.  \code{raw\_rot[j]} encodes the cumulative heading (in $30°$ units)
at the end of edge $j$ in the canonical, unreflected orientation.

% -----------------------------------------------------------------------
\section{Hat Geometry Functions}

\subsection{\code{hv14} — Local 14-Vertex Polygon}

\begin{tcolorbox}[dualbox, title={\code{hv14(d, ref, pt)} — build the 14 vertices}]
  Given a discrete direction \code{d}, a reflection flag \code{ref},
  and a pivot vertex index \code{pt}, \code{hv14} computes the 14 vertices of
  the hat in a local coordinate frame with the pivot at the origin.

  The loop over $j = 0\ldots 13$ reads each edge's properties from the two
  lookup arrays, applying a cyclic shift and optional reversal for reflected
  tiles.  For a reflected tile the traversal order is reversed
  (\code{(pt+13-j) \% 14}) and all rotations are negated
  (\code{(-rot) \% 12}).

  The heading $\alpha_j$ at step $j$ is
  \[
    \alpha_j = \pi\,(\mathit{erot}[j] + 1)/6
  \]
  so that steps of $1$ in \code{erot} correspond to $30°$ turns.
  Vertex positions are accumulated by integrating these headings.

  The final reordering \code{(i-pt) \% 14} shifts the vertex array so
  that index $0$ in the output corresponds to the pivot vertex.
\tcblower
\begin{minted}{python}
R3 = math.sqrt(3.0)

def hv14(d, ref, pt):
    elen = [0.0]*14; erot = [0]*14
    for j in range(14):
        src = (pt+(13-j))%14 if ref \
              else (pt+j)%14
        elen[j] = R3 if raw_len[src] else 1.0
        rot = raw_rot[src] + d
        erot[j] = (-rot)%12 if ref \
                  else rot%12
    vx=[0.0]*15; vy=[0.0]*15
    for j in range(14):
        a = math.pi*(erot[j]+1)/6.0
        vx[j+1]=vx[j]+elen[j]*math.cos(a)
        vy[j+1]=vy[j]+elen[j]*math.sin(a)
    return np.array(
      [(vx[(i-pt)%14], vy[(i-pt)%14])
       for i in range(14)])
\end{minted}
\end{tcolorbox}

\subsection{\code{real\_world} — CSV Parameters to World Coordinates}

\begin{tcolorbox}[dualbox, title={\code{real\_world(d, ref, piv, x, y)}}]
  Converts the five CSV parameters into an array of 14 world-space vertices.

  The trick is that \code{hv14} already places the pivot vertex at the origin
  of the local frame.  \code{real\_world} therefore:
  \begin{enumerate}
    \item Calls \code{hv14} with direction $\pm 1$ (the ``natural'' direction
      for the reflected/non-reflected case) and pivot \code{piv = (pt-1)\%14}.
    \item Builds the $2\times 2$ rotation matrix $R$ for angle
      $\theta = \mathit{dir}\cdot\pi/6$ (negated for reflected tiles).
    \item Returns $V \cdot R^{\top} + (x, y)$: rotate the local-frame
      vertices and translate to the anchor.
  \end{enumerate}
  The result is an $[\,14\times 2\,]$ array of $(x,y)$ world positions.
\tcblower
\begin{minted}{python}
def real_world(d, ref, piv, x, y):
    V = hv14(1 if ref else -1, ref, piv)
    th = d*(math.pi/6.0)*(-1 if ref else 1)
    c = math.cos(th); s = math.sin(th)
    Rm = np.array([[c,-s],[s,c]])
    return V@Rm.T + np.array([x, y])
\end{minted}
\end{tcolorbox}

\subsection{\code{cworld} — Tetrille Integer Coordinates to World Space}

The grow-overlay data is stored as tetrille integer triples $(v, x, y)$ where
$v \in \{3, 4, 6\}$ is the vertex valence.  \code{cworld} converts to
Euclidean $\mathbb{R}^2$:

\begin{minted}{python}
def cworld(v, x, y):          # tetrille Coord -> world
    if v == 6:   sx, sy = 6*x, 6*y
    elif v == 4: sx, sy = 3*x, 3*y
    else:        sx, sy = 2*(x-y), 2*(x+2*y)
    return ((3-sx-sy)/2.0, (sy-sx+3)/(2*R3))
\end{minted}

Each valence system has a different basis in terms of the underlying $v=6$
units: $v=4$ points are at half-spacing, and $v=3$ points lie on the
triangular sublattice at a further $1/\sqrt{3}$ rescaling.  The formula
converts to a shared Euclidean frame where the canonical hat tile fits
alongside the \file{hat.tile} SVG output.

% -----------------------------------------------------------------------
\section{IO Functions}

\subsection{\code{load\_csv} — Reading Hat Placement Files}

\begin{tcolorbox}[dualbox, title={\code{load\_csv(fn)} — load hat records from CSV}]
  Reads a CSV file whose data rows are of the form
  \texttt{x,y,dir,ref,pt}.  Lines are skipped if they do not begin with
  a digit, sign, or decimal point — this quietly ignores the header line and
  any comment lines.

  Each accepted row is parsed into a Python dict with keys
  \code{x}, \code{y}, \code{dir}, \code{ref}, \code{pt}, and a pre-computed
  \code{world} array ($14\times 2$ float array) obtained by calling
  \code{real\_world(d, ref, (pt-1)\%14, x, y)}.

  The $(pt-1)\%14$ shift converts the 1-based CSV pivot index to the
  0-based index expected by \code{hv14}.
\tcblower
\begin{minted}{python}
def load_csv(fn):
    out = []
    for ln in open(fn):
        s = ln.strip()
        if not s or s[0] not in \
                "-+.0123456789":
            continue
        p = s.split(",")
        x,y = float(p[0]),float(p[1])
        d = int(p[2])
        ref = 1 if p[3][0] in "Tt1" else 0
        pt = int(p[4])
        out.append(dict(
          x=x, y=y, dir=d, ref=ref, pt=pt,
          world=real_world(
            d, ref, (pt-1)%14, x, y)))
    return out
\end{minted}
\end{tcolorbox}

\subsection{\code{load\_records} — Reading Grow-Symmetric2 Output}

\begin{tcolorbox}[dualbox, title={\code{load\_records(level)} — parse .dat completion files}]
  Reads \file{data/sym12/sym\_completions\_level\{L\}.dat}, which is
  the output of \code{grow\_symmetric2}.  Each record in that file contains
  a \texttt{tiles:} field whose value is a list of hat cycles, each cycle
  being a list of tetrille integer triples.

  The function uses two nested regex passes:
  \begin{enumerate}
    \item \verb|tiles:(\[\[.*?\]\])| — extracts the bracketed tiles string for
      each record.
    \item \verb|\[((?:\(-?\d+,-?\d+,-?\d+\),?)+)\]| — splits into individual
      hat cycle strings.
    \item \verb|\((-?\d+),(-?\d+),(-?\d+)\)| — extracts each vertex triple
      $(v, x, y)$.
  \end{enumerate}
  Each triple is converted to world space via \code{cworld}, yielding a list
  of $[\,14\times 2\,]$ float arrays — one per hat tile per record.
\tcblower
\begin{minted}{python}
def load_records(level):
    txt = open(
      f"data/sym12/"
      f"sym_completions_level{level}.dat"
    ).read()
    recs = []
    for line in re.findall(
            r'tiles:(\[\[.*?\]\])', txt):
        hats = [
          np.array([cworld(*map(int,m))
            for m in re.findall(
              r'\((-?\d+),(-?\d+),(-?\d+)\)',
              hat)])
          for hat in re.findall(
            r'\[((?:\(-?\d+,-?\d+,-?\d+\)'
            r',?)+)\]', line)]
        recs.append(hats)
    return recs
\end{minted}
\end{tcolorbox}

% -----------------------------------------------------------------------
\section{Super Cluster Construction}

\subsection{Rotation Constants}

Three constants control the assembly of the super cluster:

\begin{minted}{python}
C13         = np.array((-68.5, 0.8660254037844386))  # arrangement #13 centre
COPY_ANGLES = [180.0, 300.0, 60.0]                   # 0/120/240 + 180 deg display rotation
\end{minted}

\noindent $C_{13}$ is the world-space position of the $v[12]$-type triple vertex
for arrangement \#13 — the gap-free H8 centre.  The three copy angles encode a
$120°$-step rotation ($0°/120°/240°$) combined with a global $+180°$ display
rotation that corrects for the orientation convention used by \code{rotsym.c}.

\subsection{\code{rotate\_param} — Rotating a CSV Hat Record}

\begin{tcolorbox}[dualbox, title={\code{rotate\_param(h, deg, ctr)} — rotate placement parameters}]
  Rotating a hat tile in \emph{world space} by angle \code{deg} about a
  centre \code{ctr} must be reflected in the CSV parameters $(x, y, \mathit{dir})$.

  The \code{dir} field counts $30°$ steps, so the rotation step count is
  $u = \mathit{deg}/30$.  The script asserts this is an integer.  For a
  regular hat the direction increases CCW ($+u$); for a reflected hat the
  direction increases CW ($-u$), because the chirality flip reverses the
  angular convention.

  The pivot position $(x, y)$ is rotated geometrically:
  $(nx, ny) = \mathit{ctr} + R_{\mathit{deg}}\cdot((x,y) - \mathit{ctr})$.

  A geometric cross-check (not in the production loop but run during
  verification) confirms that \code{r["world"]} matches the directly-rotated
  world array to within \textasciitilde$10^{-13}$.
\tcblower
\begin{minted}{python}
def rotate_param(h, deg, ctr):
    u = deg / 30.0
    assert abs(u-round(u)) < 1e-9
    u = int(round(u))
    nd = (h["dir"] +
          (-u if h["ref"] else u)) % 12
    nx, ny = ctr + rotmat(deg) @ (
        np.array([h["x"],h["y"]]) - ctr)
    return dict(
      x=float(nx), y=float(ny),
      dir=nd, ref=h["ref"], pt=h["pt"],
      world=real_world(
        nd, h["ref"],
        (h["pt"]-1)%14, nx, ny))
\end{minted}
\end{tcolorbox}

\subsection{Assembling the Three Copies}

The base tile set is loaded from \file{from\_python/SUPER\_SUPER\_H8\_TILE.dat}.
Each of the three copy angles is applied to every hat in the base:

\begin{minted}{python}
base = load_csv("from_python/SUPER_SUPER_H8_TILE.dat")
copies = []
for ci, ang in enumerate(COPY_ANGLES):
    copy = []
    for h in base:
        r = rotate_param(h, ang, C13)
        # geometric cross-check
        ref_world = (h["world"]-C13) @ rotmat(ang).T + C13
        maxerr = max(maxerr, np.abs(r["world"]-ref_world).max())
        r["copy"] = ci
        copy.append(r)
    copies.append(copy)
super_hats = [h for cp in copies for h in cp]
\end{minted}

\noindent The cross-check verifies that the parameter-level rotation matches the
directly-rotated world vertices to numerical precision.  The flat list
\code{super\_hats} then contains all three copies concatenated; its length is
$3 \times |\mathit{base}|$.  The centroid array
\code{Scent = np.array([h["world"].mean(0) for h in super\_hats])} stores
the 2D centre of mass of every hat's 14 vertices, used as the target point set
for the alignment step.

% -----------------------------------------------------------------------
\section{Grow Overlay: Alignment Algorithm}

\subsection{Overview}

The grow-symmetric2 completions live in a coordinate frame centred at
\[
  g_{\mathit{ctr}} = \mathit{cworld}(3, 1, -1)
\]
which is the tetrille world-space position of $v[12]$ in the canonical hat
tile.  This is \emph{not} the same as $C_{13}$; a rigid-body transform
(rotation and optional reflection) must be found to map grow-space centroids
onto super-cluster centroids.

The function \code{fullalign} performs this alignment using a procedure
analogous to the Iterative Closest Point (ICP) algorithm with a global
initialisation sweep.

\subsection{\code{fullalign} — ICP-Style Centroid Alignment}

\begin{tcolorbox}[dualbox, title={\code{fullalign(hats)} — find best rigid transform}]
  \textbf{Input:} a list of $[\,14\times 2\,]$ world arrays from one grown
  completion record.  Centroids $G$ are computed by averaging each hat's
  14 vertices.

  \medskip
  \textbf{Outer loop (reflection)} iterates over \code{refl} $\in \{+1, -1\}$.
  When \code{refl = -1} the $x$-coordinates of the zero-centred point set $P$
  are negated before any rotation, implementing a mirror.

  \medskip
  \textbf{Middle loop (initial angle)} sweeps the starting rotation in $3°$
  steps from $0°$ to $360°$.

  \medskip
  \textbf{Inner loop (ICP refinement):} up to 6 SVD-based iterations.
  At each iteration:
  \begin{enumerate}
    \item Rotate $P$ by the current estimate $R$ and add $C_{13}$ to get
      candidate world positions $T$.
    \item For each point in $T$ find its nearest neighbour in \code{Scent}
      (squared distances).
    \item Build the cross-covariance matrix $H = A^{\top}B$ where $A$ are the
      matched source points and $B$ the corresponding target points.
    \item Solve $H = U\Sigma V^{\top}$ (SVD) and set $R_k = V^{\top}U^{\top}$
      (orthogonal Procrustes).
    \item Extract the updated angle from $R_k$.
  \end{enumerate}
  The final coverage fraction $f$ (share of grown points landing within
  distance $\sqrt{6\times 10^{-2}}$ of a super-cluster centroid) is used
  to rank candidates.  The function returns the best $(f, (\mathit{refl}, R))$.
\tcblower
\begin{minted}{python}
gctr = np.array(cworld(3, 1, -1))

def fullalign(hats):
    G = np.array([h.mean(0) for h in hats])
    best = (-1, None)
    for refl in (1, -1):
        P = (G-gctr).copy()
        P[:, 0] *= refl
        for d0 in range(0, 360, 3):
            deg = float(d0)
            for _ in range(6):
                Rm = rotmat(deg)
                T = P@Rm.T + C13
                d2 = ((T[:,None]
                       -Scent[None])**2
                      ).sum(-1)
                nn = d2.argmin(1)
                msk = d2.min(1) < 6e-2
                if msk.sum() < 3: break
                A = P[msk]
                B = Scent[nn[msk]] - C13
                H = A.T@B
                U,_,Vt = np.linalg.svd(H)
                Rk = Vt.T@U.T
                deg = math.degrees(
                  math.atan2(Rk[1,0],Rk[0,0]))
            Rm = rotmat(deg)
            T = P@Rm.T + C13
            f = (((T[:,None]-Scent[None])**2
                 ).sum(-1).min(1)<6e-2).mean()
            if f > best[0]:
                best = (f, (refl, Rm))
    return best
\end{minted}
\end{tcolorbox}

\noindent The threshold $6\times 10^{-2}$ is chosen to be smaller than the
minimum inter-hat distance in world space, ensuring that each match is
unambiguous.

\subsection{Locking the Transform from Level-5}

The global rigid-body transform is fixed once, using the level-5 completion
record that achieves the best coverage:

\begin{minted}{python}
best5 = (-1, None, None)
for i, rec in enumerate(load_records(5)):
    f, tr = fullalign(rec)
    if f > best5[0]: best5 = (f, i, tr)
f5, idx5, (refl, Rm) = best5
\end{minted}

\noindent All subsequent per-level lookups use the same \code{(refl, Rm)} pair via
the closure \code{place(hats)}, which applies the fixed reflection + rotation
and translates to $C_{13}$:

\begin{minted}{python}
def place(hats):
    out = []
    for h in hats:
        p = (h-gctr).copy(); p[:, 0] *= refl
        out.append(p@Rm.T + C13)
    return out
\end{minted}

% -----------------------------------------------------------------------
\section{Per-Level Matching and Birth Assignment}

Once the transform is fixed, the script iterates over levels $L = 1\ldots 5$.
For each level it finds the completion record that covers the most of
\code{Scent} and records which super-cluster hat first appears at that level:

\begin{minted}{python}
birth = {}                   # super-hat index -> first level
for L in range(1, 6):
    bestrec = (-1, None, None)
    for i, rec in enumerate(load_records(L)):
        placed = place(rec)
        pc = np.array([h.mean(0) for h in placed])
        cov = ((pc[:,None]-Scent[None])**2).sum(-1).min(1) < 6e-2
        if cov.mean() > bestrec[0]:
            bestrec = (cov.mean(), i, placed)
    frac, ri, placed = bestrec
    pc = np.array([h.mean(0) for h in placed])
    nn = ((pc[:,None]-Scent[None])**2).sum(-1)
    for j in range(len(placed)):
        si = nn[j].argmin()
        if nn[j, si] < 6e-2 and si not in birth:
            birth[si] = L
\end{minted}

\noindent After the loop, \code{birth[si]} gives the \emph{birth level} of
super-cluster hat \code{si}: the earliest ring at which that hat is
covered by a symmetric grown completion.  This assignment drives the
level-coloured overlay in the SVG.

% -----------------------------------------------------------------------
\section{SVG Rendering}

\subsection{Coordinate Scaling}

The script computes a uniform pixel scale \code{SC} such that the widest
dimension of the super cluster spans exactly 1400 pixels, with a margin
\code{M = 24} on each side:

\begin{minted}{python}
allw  = np.vstack([h["world"] for h in super_hats])
mnx, mny = allw.min(0);  mxx, mxy = allw.max(0)
W, H = mxx-mnx, mxy-mny;  span = max(W, H)
SC = 1400 / span
PW, PH = W*SC+2*M, H*SC+2*M
def px(p): return (p[0]-mnx)*SC+M
def py(p): return PH-((p[1]-mny)*SC+M)   # SVG y flipped
\end{minted}

The $y$-axis is flipped (\code{py}) because SVG coordinates increase downward
while mathematical coordinates increase upward.

\subsection{Drawing Layers}

The SVG is built in four layers (painter's order, back to front):

\begin{enumerate}
  \item \textbf{Super cluster tiles.}  Every hat in \code{super\_hats} is drawn
    as a closed 14-vertex polygon:
    \begin{itemize}
      \item Anti-hats (\code{ref=1}): fill \texttt{\#f4d03f}
        ({\color{antiyellow}\rule{0.8em}{0.8em}} golden yellow).
      \item Regular hats (\code{ref=0}): fill \texttt{\#c5d6ea}
        ({\color{hatblue}\rule{0.8em}{0.8em}} pale steel blue).
      \item Stroke: \texttt{\#7a8aaa} at width 0.3 px.
    \end{itemize}
  \item \textbf{Level-coloured overlay.}  For each super-hat index \code{si}
    with a known birth level $L$, the same polygon is redrawn with fill colour
    from a 5-step magma-ish palette and 50\% opacity:
    \begin{center}
    \begin{tabular}{@{}cll@{}}
    \toprule
    \textbf{Level} & \textbf{Hex} & \textbf{Colour} \\
    \midrule
    1 & \texttt{\#3b0f70} & {\color{lvone}\rule{1em}{0.8em}} deep purple \\
    2 & \texttt{\#641a80} & {\color{lvtwo}\rule{1em}{0.8em}} violet \\
    3 & \texttt{\#a52c60} & {\color{lvthree}\rule{1em}{0.8em}} magenta \\
    4 & \texttt{\#de4968} & {\color{lvfour}\rule{1em}{0.8em}} coral \\
    5 & \texttt{\#fe9f6d} & {\color{lvfive}\rule{1em}{0.8em}} peach \\
    \bottomrule
    \end{tabular}
    \end{center}
    The 50\% opacity lets the anti-hat yellow show through where anti-hats
    receive an overlay.
  \item \textbf{Boundary.}  The outer boundary of the grown cluster
    (\code{birth.keys()}) is computed by \code{boundary\_edges} and drawn
    as bold black lines (width 2.2 px).
  \item \textbf{Legend and centre marker.}  A small colour legend for levels
    1–5 and the anti-hat colour.  A red circle of radius 6 px marks $C_{13}$.
\end{enumerate}

\subsection{\code{boundary\_edges} — Outer Boundary of a Tile Subset}

\begin{tcolorbox}[dualbox, title={\code{boundary\_edges(idxs)} — extract the boundary}]
  A boundary edge of a planar tile union is an edge that is shared by exactly
  one tile.  For each edge in the 14-vertex boundary of each selected tile,
  the function records it in a canonical form (sorted endpoint key pair) and
  counts how many tiles share it.  Edges with count $= 1$ are boundary edges.

  Vertex coordinates are rounded to $10^{-4}$ before hashing to avoid
  floating-point near-misses from independent computations of the same
  geometric point.
\tcblower
\begin{minted}{python}
def boundary_edges(idxs):
    cnt = defaultdict(int); seg = {}
    def k(p):
        return (round(p[0]*1e4),
                round(p[1]*1e4))
    for si in idxs:
        h = super_hats[si]["world"]
        for i in range(14):
            a=h[i]; b=h[(i+1)%14]
            ka=k(a); kb=k(b)
            e=(ka,kb) if ka<=kb \
              else (kb,ka)
            cnt[e]+=1; seg[e]=(a,b)
    return [seg[e]
            for e,c in cnt.items()
            if c == 1]
\end{minted}
\end{tcolorbox}

% -----------------------------------------------------------------------
\section{CSV Export}

\subsection{File Format}

Both output CSV files share the same header and row format:

\begin{minted}{text}
x,y,dir,ref,pt
-68.500000000000,0.866025403784,7,False,1
...
\end{minted}

\noindent Coordinates are written with 12 significant figures via
\texttt{\%g} formatting; \code{ref} is written as \texttt{True}/\texttt{False}.

\subsection{SUPER\_SUPER\_H8\_ROT.csv}

Contains all hats from \code{super\_hats} (three copies concatenated in the
order copy~0, copy~1, copy~2).  This is the input consumed by
\code{rotsym.c} and downstream C programs that need the full 3-fold cluster
in CSV form.

\subsection{LEVEL\_ROT.csv}

Contains only the hats covered by the grow overlay (those in \code{birth}),
sorted inside-out: first by birth level, then by Euclidean distance from
$C_{13}$ within each level:

\begin{minted}{python}
order = sorted(birth.keys(),
               key=lambda si: (birth[si],
                               np.hypot(*(Scent[si]-C13))))
write_csv("from_python/LEVEL_ROT.csv",
          [super_hats[si] for si in order])
\end{minted}

\noindent This ordering makes it natural to stream the file and animate
ring-by-ring growth: reading the file from top to bottom adds tiles
from the centre outward.

% -----------------------------------------------------------------------
\section{Architecture Summary}

\begin{center}
\begin{tcolorbox}[
  colframe=rulecolor, colback=white, boxrule=0.6pt, arc=3pt,
  width=0.98\linewidth
]
\begin{minted}{text}
constants: C13=(-68.5, 0.866), COPY_ANGLES=[180,300,60]

1. SUPER CLUSTER
   load_csv("from_python/SUPER_SUPER_H8_TILE.dat")  -> base (N hats)
   for ci, ang in enumerate(COPY_ANGLES):
       for h in base: rotate_param(h, ang, C13)     -> copies[ci]
   super_hats = copies[0] + copies[1] + copies[2]   (3N hats)
   Scent = centroid of each hat's 14 vertices        (3N x 2)

2. GROW OVERLAY
   gctr = cworld(3, 1, -1)   (grow-space origin)
   fullalign(rec) for each rec in load_records(5):
       ICP: refl in {+1,-1}, sweep 0..360 deg (3 deg steps),
            6 SVD-Procrustes iterations each
       -> best (fraction, (refl, Rm))
   lock transform: place(hats) = (hats-gctr)*refl @ Rm.T + C13

   for L in 1..5:
       for rec in load_records(L):
           placed = place(rec); cov = nearest-in-Scent < 6e-2
       best rec at L -> bestrec
       assign birth[si] = L for newly matched super-hat indices

3. SVG
   boundary_edges(birth.keys()) -> outer boundary segments
   draw super_hats:  ref=1 -> #f4d03f, ref=0 -> #c5d6ea
   draw overlay:     birth[si]=L -> LV[L-1] at opacity 0.5
   draw boundary:    black, stroke-width=2.2
   write from_python/h8_rot_overlay.svg

4. CSV EXPORT
   write_csv(SUPER_SUPER_H8_ROT.csv, super_hats)
   order = sorted by (birth[si], |Scent[si]-C13|)
   write_csv(LEVEL_ROT.csv, [super_hats[si] for si in order])
\end{minted}
\end{tcolorbox}
\end{center}

% -----------------------------------------------------------------------
\section{Results}

Running \code{python3 MyHatTileStuff/h8\_rotcluster.py} from the project root
prints:

\begin{minted}{text}
building #13 super cluster (centre -68.5000,0.8660, +180 deg)...
  78 hats (6 anti-hats); CSV-param vs geometry max error = 2.22e-16
  grow alignment locked from level5 rec1 = 100.0%
  level 1: rec 1 matches interior at 15.4% (cumulative covered hats=12)
  level 2: rec 1 matches interior at 30.8% (cumulative covered hats=24)
  level 3: rec 1 matches interior at 57.7% (cumulative covered hats=45)
  level 4: rec 1 matches interior at 76.9% (cumulative covered hats=60)
  level 5: rec 1 matches interior at 100.0% (cumulative covered hats=78)
wrote from_python/h8_rot_overlay.svg
wrote from_python/SUPER_SUPER_H8_ROT.csv (78 hats: 26+26+26)
wrote from_python/LEVEL_ROT.csv (78 hats; per-level {1: 12, 2: 12, 3: 21, 4: 15, 5: 18})
\end{minted}

\medskip
Key observations:
\begin{itemize}
  \item The geometric cross-check error of $2.22\times 10^{-16}$ (machine
    epsilon) confirms that the parameter-level rotation in \code{rotate\_param}
    is exactly consistent with the direct geometric rotation.
  \item The level-5 alignment achieves 100\% coverage, confirming that the
    grow-symmetric2 output at level 5 contains a completion that fills the
    entire H8 super cluster interior.
  \item The 78 hats split into 26 per copy (three-fold symmetry) and comprise
    6 anti-hats total ($2$ per copy), consistent with the expected $3:1$
    hat-to-anti-hat ratio.
\end{itemize}

% -----------------------------------------------------------------------
\section{Usage}

\begin{minted}{bash}
# Run from the project root (no arguments needed)
python3 MyHatTileStuff/h8_rotcluster.py

# Prerequisites
#   from_python/SUPER_SUPER_H8_TILE.dat   -- base H8 super tile (CSV)
#   data/sym12/sym_completions_level{1-5}.dat  -- grow_symmetric2 output

# Outputs
#   from_python/h8_rot_overlay.svg    -- rendered image
#   from_python/SUPER_SUPER_H8_ROT.csv
#   from_python/LEVEL_ROT.csv
\end{minted}

\medskip
\noindent The script requires only the Python standard library and NumPy.
No SciPy or other geometry libraries are needed.

\vspace{4pt}
\noindent\rule{\linewidth}{0.6pt}
\vspace{2pt}

{\small\color{gray}
\file{MyHatTileStuff/h8\_rotcluster.py} — part of PolyformTown.
Compile this document with \texttt{pdflatex -shell-escape h8\_rotcluster.tex}
(requires \texttt{pygments} for \texttt{minted}).
}

\end{document}
